% Adaptación de FORMATOAPE.tex para APE2; original preservado.

\documentclass[11pt,a4paper]{article}
\usepackage[utf8]{inputenc}
\usepackage[T1]{fontenc}
\usepackage[spanish]{babel}
\usepackage{geometry}
\usepackage{graphicx}
\usepackage{array}
\usepackage{tabularx}
\usepackage{makecell}
\usepackage{hyperref}
\usepackage{xcolor}
\usepackage{fancyhdr}
\usepackage{titlesec}
\usepackage{helvet}
\usepackage{ragged2e}
\usepackage{enumitem}
\usepackage{setspace}
\usepackage{amsmath}
\usepackage{booktabs}
\usepackage{float}
\usepackage{listings}
\usepackage{hyperref}
\usepackage{xurl}
\usepackage{microtype}
\usepackage{float}
\emergencystretch=2em
\renewcommand{\familydefault}{\sfdefault}
\geometry{a4paper,left=3cm,right=3cm,top=2.35cm,bottom=2.35cm,headheight=1.9cm,headsep=0.45cm}
\setlength{\parindent}{0pt}
\setlength{\parskip}{0.3em}
\setlength{\tabcolsep}{5pt}
\renewcommand{\arraystretch}{1.13}
\pagestyle{fancy}
\fancyhf{}
\fancyhead[L]{\IfFileExists{reporte_tecnico_overleaf_logo.png}{\includegraphics[height=1.35cm]{reporte_tecnico_overleaf_logo.png}}{\small Universidad Nacional de Loja}}
\fancyhead[R]{\small FEIRNNR -- Carrera de Computación}
\fancyfoot[C]{\thepage}
\renewcommand{\headrulewidth}{0.5pt}
\renewcommand{\footrulewidth}{0pt}
\titleformat{\section}{\bfseries\fontsize{15}{17}\selectfont}{\thesection.}{0.8em}{}
\titlespacing*{\section}{0pt}{1.0ex}{0.3ex}
\newcolumntype{L}[1]{>{\RaggedRight\arraybackslash}p{#1}}
\newcolumntype{Y}{>{\RaggedRight\arraybackslash}X}
\lstset{basicstyle=\ttfamily\footnotesize,breaklines=true,frame=single,showstringspaces=false,columns=fullflexible}
\hypersetup{hidelinks}

% Soporte mínimo del contenido académico extenso.
\usepackage{longtable}
\usepackage{calc}
\usepackage{fvextra}
\providecommand{\tightlist}{\setlength{\itemsep}{0pt}\setlength{\parskip}{0pt}}

\begin{document}
\thispagestyle{fancy}
\begin{center}
{\bfseries\fontsize{21}{24}\selectfont Reporte Técnico de Actividades\\ Práctico-Experimentales Nro. 002\par}
\vspace{0.25cm}
{\large Elaboración y refinamiento de specification.md con historias de usuario, requisitos, criterios de aceptación y contratos preliminares.}
\end{center}

\section{Datos de Identificación del Estudiante y la Práctica}
\begin{tabularx}{\textwidth}{|L{0.43\textwidth}|Y|}
\hline
\textbf{Nombre del estudiante(s)} & [PENDIENTE: integrantes reales del equipo] \\ \hline
\textbf{Asignatura} & Desarrollo Basado en Plataformas \\ \hline
\textbf{Ciclo} & 5 A \\ \hline
\textbf{Unidad} & 1 \\ \hline
\textbf{Resultado de aprendizaje} & Discute cómo los estándares Web influyen en el desarrollo de software, bajo los principios de solidaridad, transparencia, responsabilidad y honestidad. \\ \hline
\textbf{Práctica Nro.} & 002 \\ \hline
\textbf{Título de la Práctica} & Elaboración y refinamiento de specification.md con historias de usuario, requisitos, criterios de aceptación y contratos preliminares. \\ \hline
\textbf{Nombre del Docente} & Edison Leonardo Coronel Romero \\ \hline
\textbf{Fecha} & Semana 2; [PENDIENTE: fecha real de la sesión] \\ \hline
\textbf{Horario} & Por definir según horario académico \\ \hline
\textbf{Lugar} & Aula / laboratorio asignado \\ \hline
\textbf{Tiempo planificado en el Sílabo} & 3 horas \\ \hline
\end{tabularx}

\textbf{Proyecto:} SIGEAC / SEC-UNL.\par
\textbf{Estado:} desarrollo documental preparado; validación docente, participación y publicación pendientes de evidencia.\par
Los datos corresponden a la guía actual de Desarrollo Basado en Plataformas. ACD2 se utiliza como insumo y no se cambia retrospectivamente su asignatura.
\section{Objetivo(s) de la Práctica}\label{objetivos-de-la-pruxe1ctica}

\begin{itemize}
\tightlist
\item
  Construir y refinar la especificación funcional del proyecto mediante
  Spec-Driven Development (SDD).
\item
  Transformar las necesidades de los actores en requisitos funcionales y
  no funcionales claros y verificables.
\item
  Definir historias de usuario y criterios de aceptación que permitan
  comprobar el comportamiento esperado.
\item
  Establecer contratos preliminares que sirvan de base para el diseño
  posterior de la API.
\item
  Validar la especificación con el docente e incorporar la
  retroalimentación de manera trazable.
\end{itemize}

Los cuatro primeros objetivos tienen desarrollo documental en los
archivos de esta carpeta. El quinto requiere revisión real de Edison
Leonardo Coronel Romero; no se acredita mediante una redacción supuesta.

\section{Materiales, Reactivos, Equipos y
Herramientas}\label{materiales-reactivos-equipos-y-herramientas}

\subsection{Insumos solicitados y
disponibles}\label{insumos-solicitados-y-disponibles}

{\def\LTcaptype{none} % do not increment counter
\begin{longtable}[]{@{}
  >{\raggedright\arraybackslash}p{(\linewidth - 4\tabcolsep) * \real{0.3333}}
  >{\raggedright\arraybackslash}p{(\linewidth - 4\tabcolsep) * \real{0.3333}}
  >{\raggedright\arraybackslash}p{(\linewidth - 4\tabcolsep) * \real{0.3333}}@{}}
\toprule\noalign{}
\begin{minipage}[b]{\linewidth}\raggedright
Insumo de la guía
\end{minipage} & \begin{minipage}[b]{\linewidth}\raggedright
Evidencia localizada
\end{minipage} & \begin{minipage}[b]{\linewidth}\raggedright
Tratamiento
\end{minipage} \\
\midrule\noalign{}
\endhead
\bottomrule\noalign{}
\endlastfoot
Banco de proyectos y análisis ACD1 & No se aportó/localizó una entrega
académica nominal ACD1 en esta preparación. & Recuperar contexto de
\path|docs/requirements/spec.md|; no inventar el banco ni una actividad
realizada. \\
Actores y fundamentos SDD de AA1 & No se aportó/localizó AA1 como
entrega independiente. & Utilizar actores de especificaciones y
constitución; incorporar AA1 real cuando exista. \\
Requisitos y observaciones de ACD2 &
\path|documentos-academicos/ACD2-semana-2/ACD2_SEMANA_2.md|. & Usar
refinamientos técnicos; retroalimentación de aula sigue pendiente. \\
Repositorio configurado APE1 & Git local y origin
\path|git@github.com:M16artz/SIGEAC.git|. & Se verifica configuración
actual, no la ejecución de APE1 por integrantes. \\
Documentación del proyecto & \path|docs/requirements/|, \path|specs/|,
\path|docs/testing/|. & Fuentes de requisitos, estados y
trazabilidad. \\
Diagnóstico previo & \path|AUDITORIA_DOCUMENTACION_SDD.md|. &
Identificar inconsistencias sin tratar el diagnóstico como
aprobación. \\
\end{longtable}
}

\subsection{Herramientas y función
real}\label{herramientas-y-funciuxf3n-real}

Se utilizaron archivos del repositorio, Git de solo lectura para
identificar rama/commit/remote, Markdown, inspección de fuentes Java y
consulta de documentación técnica. La estructura de historias,
requisitos, entidades y criterios se tomó de
\path|.specify/templates/spec-template.md|. No se ejecutó el CLI de Spec
Kit ni se crearon ramas, tareas de implementación o migraciones por una
actividad documental.

Quarto/Pandoc se utiliza para convertir el contenido a LaTeX, respetando
la plantilla APE original. Se conservará su recurso de logo disponible.
No hay motor TeX local disponible en PATH; la compilación y revisión
visual del PDF quedan pendientes de Overleaf.

No se realizaron instalaciones, llamadas a APIs del proyecto, pruebas de
carga ni despliegues. Tampoco se atribuye uso efectivo de VS Code o
herramientas colaborativas a estudiantes sin evidencia.

\section{Procedimiento / Metodología
Ejecutada}\label{procedimiento-metodologuxeda-ejecutada}

\subsection{Revisar los insumos
existentes}\label{revisar-los-insumos-existentes}

Se revisaron contexto, actores, alcance, exclusiones y decisiones CL de
\path|docs/requirements/spec.md|, reglas de dominio, catálogo P01--P16,
especificación 005 y ACD2. El código permite contrastar catálogo, sesión
y preparación de proceso. Los resultados de pruebas anteriores no se
presentan como ejecución nueva.

Hallazgos que orientaron el refinamiento:

\begin{itemize}
\tightlist
\item
  El problema funcional está descrito, pero falta evidencia empírica del
  proceso institucional actual.
\item
  La denominación SIGEAC/SEC-UNL debe explicitarse para identificar el
  mismo repositorio.
\item
  Las competencias institucionales y vigencia normativa siguen
  pendientes.
\item
  Los documentos históricos pueden llamar borrador al estado que el
  código actual crea como PREPARATION.
\item
  Crear proceso y activarlo son operaciones distintas.
\item
  RNF-003 dice fallos consecutivos, mientras el repositorio cuenta
  fallos en ventana; no se escoge una política sin decisión.
\item
  El OpenAPI común contiene diferencias con controladores; contratos
  preliminares se contrastaron con código.
\end{itemize}

\subsection{Refinar problema y
alcance}\label{refinar-problema-y-alcance}

Se propuso un incremento U1 de acceso autorizado, catálogo de carreras y
preparación de proceso con modelo publicado. El objetivo es identificar
carrera, período, versión y actor antes de decisiones institucionales
posteriores.

Se excluyeron del recorte autorización/activación, planificación
avanzada, evidencias, revisión, valoración, informes, mejora y archivo.
La exclusión corresponde a la selección de esta práctica; varios módulos
existen en el proyecto. La propuesta es viable como trabajo de
especificación en 3 horas, no como promesa de implementar todos los
comportamientos en ese tiempo.

La validación del recorte se registra como CL-APE-01 pendiente. No se
introducen capacidades nuevas ajenas al código o a la documentación.

\subsection{Formular requisitos
funcionales}\label{formular-requisitos-funcionales}

Se redactaron diez requisitos locales APE-FR-01--10 en
\url{specification.md}, vinculados a FR canónicos. Se separan alta de
carrera, consulta, modificación histórica, selección de modelo y
preparación de proceso. Se describen actor, condición y resultado
observable sin imponer una clase, tabla o framework como requisito de
negocio.

Ejemplo de refinamiento: ``gestionar carreras'' se desglosa en registrar
datos válidos y código único; consultar catálogo autorizado;
modificar/desactivar conservando identidad e historial. Cada operación
tiene escenarios de éxito y rechazo propios.

\subsection{Formular requisitos no
funcionales}\label{formular-requisitos-no-funcionales}

Se seleccionaron diez verificaciones APE-NFR-01--10 relativas a
seguridad, rendimiento, disponibilidad, concurrencia, auditoría,
accesibilidad, mantenibilidad, interoperabilidad, usabilidad y datos
sintéticos. Se mantienen los umbrales existentes de sesión y p95; el
entorno experimental y la aplicación de usabilidad al recorte requieren
confirmación.

Se distingue una condición verificable de una certificación: disponer de
ETag no prueba toda la concurrencia; disponer de etiquetas no acredita
WCAG AA; usar JSON no demuestra por sí solo interoperabilidad. No se
inventa un porcentaje de disponibilidad ni una medición de rendimiento.

\subsection{Elaborar historias de
usuario}\label{elaborar-historias-de-usuario}

Se elaboraron tres recorridos priorizados:

{\def\LTcaptype{none} % do not increment counter
\begin{longtable}[]{@{}
  >{\raggedright\arraybackslash}p{(\linewidth - 4\tabcolsep) * \real{0.3333}}
  >{\raggedright\arraybackslash}p{(\linewidth - 4\tabcolsep) * \real{0.3333}}
  >{\raggedright\arraybackslash}p{(\linewidth - 4\tabcolsep) * \real{0.3333}}@{}}
\toprule\noalign{}
\begin{minipage}[b]{\linewidth}\raggedright
Historia
\end{minipage} & \begin{minipage}[b]{\linewidth}\raggedright
Necesidad y beneficio
\end{minipage} & \begin{minipage}[b]{\linewidth}\raggedright
Prueba independiente planteada
\end{minipage} \\
\midrule\noalign{}
\endhead
\bottomrule\noalign{}
\endlastfoot
US1, P1: acceder y consultar funciones & N-01; terminar acceso y operar
dentro de funciones vigentes. & Identidad sintética preexistente: login,
identidad y logout. \\
US2, P1: catálogo identificable e histórico & N-02; disponer de carrera
inequívoca sin borrar historial. & Administrador/unidad existentes:
alta, consulta, cambio y desactivación. \\
US3, P2: preparar proceso & N-03; identificar evaluación sin activación
institucional implícita. & Carrera/modelo publicados existentes: alta y
consulta de PREPARATION. \\
\end{longtable}
}

Independencia significa poder validar el recorrido con prerrequisitos
disponibles, no ausencia total de dependencias del dominio.

\subsection{Definir criterios de
aceptación}\label{definir-criterios-de-aceptaciuxf3n}

Se formularon 18 escenarios Dado--Cuando--Entonces: cuatro de acceso,
ocho de catálogo y seis de preparación. Cubren éxito, campos/fechas
inválidas, duplicados, ámbito, cierre/expiración de sesión y
concurrencia.

Ejemplo: Dado que la carrera/unidad están activas, el período es válido
y el modelo está publicado, cuando el coordinador con capacidad de
gestión crea un proceso, entonces queda PREPARATION con las referencias
exactas y sin activarse. Los escenarios identificados se conectan a
requisitos y posibles pruebas; no se declara que ya fueron aprobados.

\subsection{Ejecutar clarificación}\label{ejecutar-clarificaciuxf3n}

Se identificaron y registraron las ambigüedades. Se adoptaron dos
decisiones editoriales verificables: nombre \path|specification.md|
dentro de la carpeta académica, y uso de los datos de esta guía. Se
dejaron seis preguntas CL-APE pendientes sobre recorte, competencias,
intentos fallidos, problema empírico, protocolos y retroalimentación.

\textbf{No se completó una clarificación con el docente}, porque no se
proporcionaron respuestas ni una sesión verificable. Este paso queda
preparado para realizarlo honestamente, con campos de persona, fecha,
respuesta y respaldo.

\subsection{Identificar contratos
preliminares}\label{identificar-contratos-preliminares}

Se redactaron nueve fichas CP-01--09 en \url{CONTRATOS_PRELIMINARES.md}.
Relacionan necesidad/requisito, recurso, operación, entrada,
precondición, resultado y errores. Las rutas se recuperaron de
controladores reales, pero no se exige formalizar OpenAPI definitivo en
esta práctica.

Se conservaron diferencias importantes del código: entrada de carrera
plana y respuesta con \path|data|; revisión inicial de carrera 1 y de
proceso 0; errores de campos 400 frente a algunas reglas 422; catálogo
paginado y listado de procesos sin paginación declarada. La sesión
utiliza cookie y CSRF, no JWT supuesto.

\subsection{Actualizar specification.md dentro de la estructura de Spec
Kit}\label{actualizar-specification.md-dentro-de-la-estructura-de-spec-kit}

Se conservaron las secciones \path|User Scenarios & Testing|,
\path|Requirements|, \path|Key Entities|, \path|Success Criteria| y
\path|Assumptions| de la plantilla local. Se añadieron contexto,
actores, alcance, restricciones, clarificaciones, referencias a
contratos y matriz de trazabilidad.

La ubicación de entrega es
\path|documentos-academicos/APE2-semana-2/specification.md|. Se conserva
la nomenclatura \path|spec.md| de los paquetes canónicos para no romper
sus enlaces. Esta es una selección académica enlazada con el producto;
no se afirma que haya reemplazado ni aprobado la especificación global.

Se guardó una versión de trabajo inicial real de esta preparación y una
versión refinada, junto a un diff y hashes. No se presenta el borrador
inicial como trabajo de la semana anterior ni como un commit
inexistente.

\section{Resultados}\label{resultados}

\subsection{Productos preparados}\label{productos-preparados}

{\def\LTcaptype{none} % do not increment counter
\begin{longtable}[]{@{}
  >{\raggedright\arraybackslash}p{(\linewidth - 4\tabcolsep) * \real{0.3333}}
  >{\raggedright\arraybackslash}p{(\linewidth - 4\tabcolsep) * \real{0.3333}}
  >{\raggedright\arraybackslash}p{(\linewidth - 4\tabcolsep) * \real{0.3333}}@{}}
\toprule\noalign{}
\begin{minipage}[b]{\linewidth}\raggedright
Producto solicitado
\end{minipage} & \begin{minipage}[b]{\linewidth}\raggedright
Resultado local
\end{minipage} & \begin{minipage}[b]{\linewidth}\raggedright
Límite
\end{minipage} \\
\midrule\noalign{}
\endhead
\bottomrule\noalign{}
\endlastfoot
specification.md actualizado & Documento refinado con contexto, diez FR
locales, diez NFR, tres historias, 18 AC y criterios de éxito. &
Pendiente de validar recorte/políticas. \\
Contratos preliminares & Nueve fichas con requisitos, recursos y
comportamientos. & No son OpenAPI final ni tráfico capturado. \\
Historial de cambios & Borrador v01, especificación refinada, diff y
hashes. & Historial local de esta preparación; no sustituye un commit
publicado. \\
Nota individual & Desarrollo orientativo y campos de aporte/enlace. &
Identidad, revisión personal y commit/PR pendientes. \\
Enlace al repositorio & Origin deriva en
\path|https://github.com/M16artz/SIGEAC|. &  \\
\end{longtable}
}

\subsection{Correspondencia de
calidad}\label{correspondencia-de-calidad}

{\def\LTcaptype{none} % do not increment counter
\begin{longtable}[]{@{}
  >{\raggedright\arraybackslash}p{(\linewidth - 2\tabcolsep) * \real{0.5000}}
  >{\raggedright\arraybackslash}p{(\linewidth - 2\tabcolsep) * \real{0.5000}}@{}}
\toprule\noalign{}
\begin{minipage}[b]{\linewidth}\raggedright
Criterio
\end{minipage} & \begin{minipage}[b]{\linewidth}\raggedright
Evidencia
\end{minipage} \\
\midrule\noalign{}
\endhead
\bottomrule\noalign{}
\endlastfoot
Claridad & Roles, acciones y resultados diferenciados; creación no
equivale a activación. \\
Especificidad & Código normalizado, vigencia, modelo publicado, estado
inicial y revisión de cambios. \\
Pertinencia & Cada requisito se relaciona con N-01--04 y con el recorte
U1 propuesto. \\
Verificabilidad & 18 AC y comprobaciones RNF con métodos y límites
indicados. \\
Consistencia & IDs locales enlazados a FR/RNF canónicos; exclusiones
expresas y políticas pendientes. \\
Transparencia & Fuentes identificadas; sin resultados, asistencia o
retroalimentación inventados. \\
\end{longtable}
}

\subsection{Influencia de estándares Web y
valores}\label{influencia-de-estuxe1ndares-web-y-valores}

HTTP permite distinguir recurso, método, representación y resultado; los
contratos hacen visibles esas decisiones entre cliente y servicio.
OpenAPI puede formalizar posteriormente la descripción de la API, pero
un documento válido no demuestra implementación o permisos correctos.
Spec Kit organiza el trabajo guiado por especificación; el uso de su
estructura no certifica por sí solo el resultado.

La solidaridad se concreta en compartir un vocabulario y revisar
escenarios del equipo; la transparencia, en declarar fuentes y
pendientes; la responsabilidad, en exigir verificación y proteger datos;
la honestidad, en no presentar borradores, fixtures o correcciones
asistidas como validación institucional o contribución personal
demostrada.

\section{Preguntas de Control}\label{preguntas-de-control}

\subsection{¿Qué diferencia existe entre una necesidad, un requisito y
una decisión de
implementación?}\label{quuxe9-diferencia-existe-entre-una-necesidad-un-requisito-y-una-decisiuxf3n-de-implementaciuxf3n}

La necesidad expresa el problema o beneficio del actor: identificar qué
carrera y modelo corresponden a una evaluación. El requisito especifica
un comportamiento exigible: crear un proceso con carrera, período y
versión publicada en preparación. La decisión de implementación
determina cómo resolverlo: un controlador Spring que delega a un
servicio y persiste en PostgreSQL. La necesidad justifica; el requisito
permite comprobar; la decisión técnica puede cambiar conservando el
comportamiento.

\subsection{¿Qué característica permite considerar que un requisito es
verificable?}\label{quuxe9-caracteruxedstica-permite-considerar-que-un-requisito-es-verificable}

Debe permitir decidir con evidencia si se cumple: condiciones de
entrada, acción, resultado observable y criterio de aprobación/rechazo.
``Catálogo rápido'' no basta; p95 de resultado visible ≤ 2 s en un
entorno y carga definidos sí permite una medición. Si el entorno falta,
el requisito todavía necesita precisión experimental. ``Crear proceso en
PREPARATION y no activarlo'' puede comprobarse consultando el estado
guardado y ausencia de activación.

\subsection{¿Por qué una historia de usuario debe expresar el beneficio
esperado?}\label{por-quuxe9-una-historia-de-usuario-debe-expresar-el-beneficio-esperado}

El beneficio explica para qué existe la capacidad y permite priorizarla
según una necesidad. En US3 no se solicita llenar un formulario por sí
mismo: se busca identificar la autoevaluación antes de su autorización.
Esa intención permite rechazar una activación implícita aunque la
interfaz parezca más cómoda y ayuda a evitar funcionalidades ajenas al
incremento.

\subsection{¿Qué relación existe entre una historia y sus criterios de
aceptación?}\label{quuxe9-relaciuxf3n-existe-entre-una-historia-y-sus-criterios-de-aceptaciuxf3n}

La historia expresa actor, capacidad y beneficio; sus criterios
delimitan qué resultados hacen aceptable esa capacidad. Una historia
puede necesitar varios escenarios: alta válida, código repetido, unidad
inexistente y acceso sin permiso. Los criterios no describen cada línea
de código, pero son una base para diseñar pruebas. Aprobar un único caso
feliz no cubre todos los criterios de la historia.

\subsection{¿Qué debe hacerse cuando existe una ambigüedad que el equipo
aún no puede
resolver?}\label{quuxe9-debe-hacerse-cuando-existe-una-ambiguxfcedad-que-el-equipo-auxfan-no-puede-resolver}

Registrar una pregunta identificada, alternativas, impacto, responsable
de decidir y estado pendiente; consultar al docente o persona
competente. Mientras tanto, limitar el alcance o usar un supuesto
explícito reversible sin presentarlo como aprobación. Aquí CL-APE-03
conserva la diferencia entre fallos consecutivos y acumulados; no se
cambia una política de seguridad por conveniencia de redacción.

\subsection{¿Por qué los contratos definidos en esta práctica se
consideran
preliminares?}\label{por-quuxe9-los-contratos-definidos-en-esta-pruxe1ctica-se-consideran-preliminares}

Porque identifican recursos, operaciones y datos para acordar
comportamiento, pero todavía requieren validación de alcance, esquemas,
errores, permisos y compatibilidad con clientes. Que las fichas se
apoyen en código existente reduce especulación, pero no les concede
aprobación docente ni las convierte en un OpenAPI final. Podrán
refinarse en el diseño posterior conservando su vínculo con requisitos y
evitando romper contratos operativos.

\subsection{¿Cómo se relacionará posteriormente un requisito con una
prueba
automatizada?}\label{cuxf3mo-se-relacionaruxe1-posteriormente-un-requisito-con-una-prueba-automatizada}

Se enlaza necesidad → requisito → historia → criterio → prueba →
ejecución. Por ejemplo, APE-FR-03 se relaciona con AC-U2-02: crear dos
carreras con código normalizado equivalente debe rechazar la segunda y
dejar una sola. La prueba automatizada comprobará estado HTTP o
excepción de dominio y persistencia, utilizando datos sintéticos. Se
registrarán commit, entorno, fecha y resultado. La existencia de
\path|CareerAdministrationServiceTest| es un punto de partida, no una
nueva ejecución ni prueba de todo el requisito.

\section{Conclusiones}\label{conclusiones}

Se preparó una especificación académica acotada al acceso, catálogo y
preparación de procesos, con necesidades, actores y límites explícitos.
La separación de alta y activación permite conservar la intención
institucional sin inventar competencias. Los requisitos se relacionan
con historias, escenarios y operaciones preliminares, lo que facilita su
revisión y posterior automatización.

La práctica deja pendientes verificables: aceptación docente del
alcance, evidencia institucional del problema, resolución de políticas
abiertas, protocolos RNF y contribuciones individuales. La validación
docente y publicación no se sustituyen por un documento elaborado ni por
un enlace al commit base.

\section{Recomendaciones}\label{recomendaciones}

\begin{itemize}
\tightlist
\item
  Presentar \path|specification.md| al docente y completar CL-APE con
  respuestas y respaldo real.
\item
  Incorporar ACD1, AA1 y APE1 si el equipo dispone de esos archivos; no
  reconstruirlos como hechos pasados.
\item
  Validar los contratos contra controladores y corregir OpenAPI en una
  tarea posterior autorizada.
\item
  Registrar el aporte de cada integrante con revisión personal, cambio
  verificable y enlace real de commit/PR cuando se publique.
\item
  Versionar los archivos de esta actividad de forma selectiva; no
  incluir otros archivos locales ni secretos por usar una operación Git
  indiscriminada.
\item
  Compilar en Overleaf y revisar tablas/enlaces antes de entregar el PDF
  si el docente lo solicita.
\end{itemize}

\section{Bibliografía / Referencias}\label{bibliografuxeda-referencias}

\begin{itemize}
\tightlist
\item
  Guía de Actividades Práctico-Experimentales Nro. 002, Desarrollo
  Basado en Plataformas, proporcionada por el usuario; datos y preguntas
  reproducidos en este reporte.
\item
  Proyecto SIGEAC/SEC-UNL, commit
  \path|faf868ea5a8a85285290c06d3baf8fa49889edac|:
  \path|docs/requirements/spec.md|,
  \path|docs/requirements/domain-rules.md|,
  \path|docs/requirements/process-flows.md|,
  \path|specs/005-complete-evaluation-workflows/spec.md| y
  \path|.specify/templates/spec-template.md|.
\item
  \href{../ACD2-semana-2/ACD2_SEMANA_2.md}{ACD2, insumo académico} y
  \href{../../AUDITORIA_DOCUMENTACION_SDD.md}{auditoría documental}.
\item
  \href{https://github.com/github/spec-kit}{GitHub Spec Kit, repositorio
  oficial}. Se utiliza como referencia del enfoque y se inspecciona la
  plantilla local; no se declara ejecución del CLI.
\item
  \href{https://www.rfc-editor.org/rfc/rfc9110.html}{RFC 9110: HTTP
  Semantics}, Fielding, Nottingham y Reschke, 2022.
\item
  \href{https://spec.openapis.org/oas/latest.html}{OpenAPI
  Specification}, OpenAPI Initiative; referencia de descripción de
  interfaces HTTP, no actualización de la versión del proyecto.
\end{itemize}

\section{Anexos}\label{anexos}

{\def\LTcaptype{none} % do not increment counter
\begin{longtable}[]{@{}
  >{\raggedright\arraybackslash}p{(\linewidth - 2\tabcolsep) * \real{0.5000}}
  >{\raggedright\arraybackslash}p{(\linewidth - 2\tabcolsep) * \real{0.5000}}@{}}
\toprule\noalign{}
\begin{minipage}[b]{\linewidth}\raggedright
Anexo
\end{minipage} & \begin{minipage}[b]{\linewidth}\raggedright
Archivo
\end{minipage} \\
\midrule\noalign{}
\endhead
\bottomrule\noalign{}
\endlastfoot
A. Especificación refinada & \url{specification.md} \\
B. Contratos preliminares & \url{CONTRATOS_PRELIMINARES.md} \\
C. Historial local y evidencia de diferencias &
\url{HISTORIAL_CAMBIOS.md} \\
D. Nota individual para personalizar & \url{NOTA_INDIVIDUAL.md} \\
\end{longtable}
}

La versión LaTeX incorpora la especificación, los contratos y la nota
individual como anexos de contenido. El historial/diff y manifiesto
permanecen como evidencia digital enlazada, sin capturas ficticias.

\subsection{Anexo A. Especificación
refinada}\label{anexo-a.-especificaciuxf3n-refinada}

\textbf{Feature Branch}: \path|main|, rama observada; esta práctica
documental no crea una rama de implementación.

\textbf{Created}: 2026-10-31, fecha de preparación del documento; no
equivale a fecha de asistencia a clase.

\textbf{Status}: Draft académico refinado --- pendiente de validación
docente y del equipo.

\textbf{Input}: Guía APE2 Nro. 002, Desarrollo Basado en Plataformas,
Unidad 1. Insumos disponibles: especificaciones del repositorio,
auditoría documental y ACD2. Los documentos nominales ACD1, AA1 y APE1
no fueron aportados/localizados como entregas académicas independientes
en esta preparación.

\textbf{Repositorio}:
\href{https://github.com/M16artz/SIGEAC}{M16artz/SIGEAC}, dirección
derivada del remote origin; publicación de esta entrega y accesibilidad
externa no verificadas.

\textbf{Línea base inspeccionada}:
\path|faf868ea5a8a85285290c06d3baf8fa49889edac|.

\subsubsection{Contexto, problema y
necesidades}\label{contexto-problema-y-necesidades}

SIGEAC, denominado SEC-UNL en los documentos funcionales, apoya la
autoevaluación interna de carreras. El problema abordado por la
especificación es identificar qué carrera se evalúa, qué versión del
modelo se utiliza, quién puede actuar y cómo se conservan datos y
decisiones sin confundir recopilación, revisión documental y valoración
académica.

Para este incremento, la necesidad se concentra en preparar información
de partida fiable y autorizada antes de iniciar actuaciones
institucionales. No se afirma que se hayan medido pérdidas documentales,
tiempos de gestión o errores en la institución: faltan entrevistas y
diagnóstico empírico verificables.

{\def\LTcaptype{none} % do not increment counter
\begin{longtable}[]{@{}
  >{\raggedright\arraybackslash}p{(\linewidth - 4\tabcolsep) * \real{0.3333}}
  >{\raggedright\arraybackslash}p{(\linewidth - 4\tabcolsep) * \real{0.3333}}
  >{\raggedright\arraybackslash}p{(\linewidth - 4\tabcolsep) * \real{0.3333}}@{}}
\toprule\noalign{}
\begin{minipage}[b]{\linewidth}\raggedright
Necesidad identificada
\end{minipage} & \begin{minipage}[b]{\linewidth}\raggedright
Actor
\end{minipage} & \begin{minipage}[b]{\linewidth}\raggedright
Fuente existente
\end{minipage} \\
\midrule\noalign{}
\endhead
\bottomrule\noalign{}
\endlastfoot
N-01: acceder únicamente a operaciones y datos autorizados. & Usuario
habilitado y consultor autorizado. & FR-001--006 de
\path|docs/requirements/spec.md|. \\
N-02: identificar carreras, unidades y vigencias sin códigos ambiguos ni
pérdida histórica. & Administrador funcional. & FR-043 de
\path|specs/005-complete-evaluation-workflows/spec.md|; P02. \\
N-03: preparar una evaluación con carrera, período y versión publicada
exacta. & Coordinador con función autorizada para gestionar procesos. &
FR-011--012; P05; \path|docs/requirements/domain-rules.md|. \\
N-04: comprender rechazos y resolver cambios concurrentes sin
sobrescritura silenciosa. & Administrador/coordinador. & FR-039; RNF-011
y RNF-013. \\
\end{longtable}
}

\subsubsection{Actores}\label{actores}

{\def\LTcaptype{none} % do not increment counter
\begin{longtable}[]{@{}
  >{\raggedright\arraybackslash}p{(\linewidth - 4\tabcolsep) * \real{0.3333}}
  >{\raggedright\arraybackslash}p{(\linewidth - 4\tabcolsep) * \real{0.3333}}
  >{\raggedright\arraybackslash}p{(\linewidth - 4\tabcolsep) * \real{0.3333}}@{}}
\toprule\noalign{}
\begin{minipage}[b]{\linewidth}\raggedright
Actor
\end{minipage} & \begin{minipage}[b]{\linewidth}\raggedright
Responsabilidad del incremento
\end{minipage} & \begin{minipage}[b]{\linewidth}\raggedright
Límite
\end{minipage} \\
\midrule\noalign{}
\endhead
\bottomrule\noalign{}
\endlastfoot
Usuario habilitado & Iniciar/cerrar sesión y consultar su identidad y
funciones. & Tener sesión no concede administración. \\
Administrador funcional autorizado & Crear y mantener unidades/carreras
y consultar historial. & Requiere capacidad de catálogo; no se infiere
autoridad académica. \\
Coordinador autorizado & Consultar modelos disponibles y preparar un
proceso. & Para crear requiere la capacidad de gestionar procesos en la
carrera; el nombre del cargo no basta. \\
Consultor autorizado & Leer catálogo/procesos según su asignación
vigente. & No modifica recursos por permiso de lectura. \\
Administrador técnico / cuenta técnica & Operación técnica según las
políticas existentes. & No adopta decisiones institucionales; no es
actor de alta ordinaria del catálogo en el servicio actual. \\
\end{longtable}
}

\subsubsection{Objetivo del incremento y alcance para la Unidad
1}\label{objetivo-del-incremento-y-alcance-para-la-unidad-1}

Permitir especificar y comprobar un recorrido reducido: autenticar una
identidad autorizada, disponer de una carrera identificable, seleccionar
una versión publicada y crear un proceso en preparación, con errores y
cambios de catálogo trazables.

\textbf{Incluye}: sesión y consulta de identidad; listado/consulta de
unidades; alta de unidad cuando el actor esté autorizado; alta,
consulta, modificación y desactivación de carrera con historial;
consulta de versión publicada como prerrequisito; creación y consulta de
proceso en \path|PREPARATION|; controles de ámbito, concurrencia,
auditoría y errores aplicables.

\textbf{Excluye de este incremento académico}: asignación de nuevas
autoridades; publicación o edición del modelo; autorización y activación
institucional del proceso; equipos, actividades, evidencias, revisión
documental, cálculos, informes, mejora y archivo. Son exclusiones de
esta selección, no afirmaciones de ausencia en el repositorio.

\textbf{Exclusiones del producto inicial que se conservan}: acreditación
oficial, evaluación externa, integración institucional productiva,
aplicación móvil, operación sin conexión y datos institucionales reales
no autorizados.

\textbf{Viabilidad}: la práctica dura 3 horas y produce especificación,
clarificaciones y contratos preliminares; no se plantea construir ni
certificar todo el sistema en ese tiempo. El recorte propuesto debe ser
confirmado por el docente (CL-APE-01).

\subsubsection{Restricciones y separación de
decisiones}\label{restricciones-y-separaciuxf3n-de-decisiones}

\begin{itemize}
\tightlist
\item
  Mantener la separación entre revisión documental, valoración y
  acreditación, aun cuando no se ejecuten aquí.
\item
  Usar identidades y datos sintéticos; conservar los identificadores
  canónicos FR/RNF y referenciar sus refinamientos.
\item
  No cambiar políticas institucionales ni presentar datos demo como
  normativa corroborada.
\item
  Las tecnologías existentes son contexto de contraste, no parte de la
  redacción de necesidades: Angular, Spring Boot, PostgreSQL y el
  backend modular se describen fuera de los requisitos funcionales.
\item
  Conservar los archivos históricos; las correcciones académicas quedan
  en este documento y en el registro de cambios, sin sustituir
  silenciosamente las especificaciones del producto.
\end{itemize}

\subsubsection{User Scenarios \& Testing}\label{user-scenarios-testing}

Esta sección conserva la estructura de
\path|.specify/templates/spec-template.md|: historias priorizadas,
motivo, prueba independiente y escenarios de aceptación. Se utiliza
Dado--Cuando--Entonces en español. Las pruebas se describen para
validación; no se declara su ejecución en esta práctica.

\paragraph{User Story 1 - Acceder y consultar funciones vigentes
(Priority:
P1)}\label{user-story-1---acceder-y-consultar-funciones-vigentes-priority-p1}

Como usuario habilitado, quiero iniciar y cerrar mi sesión y consultar
mis funciones vigentes, para operar solo mientras mi acceso esté
autorizado y poder terminarlo cuando finalice mi trabajo.

\textbf{Why this priority}: protege el acceso a los demás recorridos y
proporciona valor independiente al comprobar identidad y ámbitos.

\textbf{Independent Test}: utilizar una identidad sintética ya
habilitada; iniciar sesión, consultar identidad, cerrar sesión e
intentar repetir la consulta, sin necesitar crear una carrera o proceso.

\textbf{Acceptance Scenarios}:

\begin{enumerate}
\def\labelenumi{\arabic{enumi}.}
\tightlist
\item
  \textbf{AC-U1-01}: \textbf{Dado que} la identidad está activa y sus
  credenciales son válidas, \textbf{cuando} inicia sesión y consulta su
  identidad, \textbf{entonces} obtiene sesión y funciones vigentes sin
  datos de otras identidades.
\item
  \textbf{AC-U1-02}: \textbf{Dado que} la sesión está activa,
  \textbf{cuando} se cierra y se intenta reutilizar, \textbf{entonces}
  la operación protegida es rechazada.
\item
  \textbf{AC-U1-03}: \textbf{Dado que} las credenciales son incorrectas
  o la identidad no existe, \textbf{cuando} intenta autenticarse,
  \textbf{entonces} no se concede sesión y el rechazo no confirma la
  existencia de la cuenta.
\item
  \textbf{AC-U1-04}: \textbf{Dado que} han transcurrido 30 minutos sin
  actividad o la identidad fue desactivada, \textbf{cuando} solicita una
  operación protegida, \textbf{entonces} debe autenticarse nuevamente y
  no recibe el recurso protegido.
\end{enumerate}

\paragraph{User Story 2 - Mantener catálogo identificable e histórico
(Priority:
P1)}\label{user-story-2---mantener-catuxe1logo-identificable-e-histuxf3rico-priority-p1}

Como administrador funcional autorizado, quiero registrar y mantener
carreras con unidad y vigencia, para preparar evaluaciones sobre un
catálogo identificable sin borrar sus relaciones históricas.

\textbf{Why this priority}: la carrera es la unidad evaluada y su
identidad condiciona el proceso posterior.

\textbf{Independent Test}: con administrador y unidad sintéticos
disponibles, crear una carrera, consultar, modificar con revisión
vigente, desactivar y comprobar historial sin crear un proceso. La alta
de unidad es un prerrequisito o alternativa administrativa, no obliga a
implementar todas las historias juntas.

\textbf{Acceptance Scenarios}:

\begin{enumerate}
\def\labelenumi{\arabic{enumi}.}
\tightlist
\item
  \textbf{AC-U2-01}: \textbf{Dado que} la unidad existe y está activa y
  los datos son válidos, \textbf{cuando} registra una carrera con código
  no utilizado tras normalización, \textbf{entonces} se crea con
  identidad, unidad, vigencia y revisión inicial.
\item
  \textbf{AC-U2-02}: \textbf{Dado que} el código normalizado ya existe,
  \textbf{cuando} intenta registrar otro equivalente, \textbf{entonces}
  se rechaza y no aparece un duplicado.
\item
  \textbf{AC-U2-03}: \textbf{Dado que} falta un campo obligatorio, la
  unidad no está disponible o la vigencia termina antes de iniciar,
  \textbf{cuando} envía el alta, \textbf{entonces} se informa el error
  correspondiente y no se crea la carrera.
\item
  \textbf{AC-U2-04}: \textbf{Dado que} consulta la revisión vigente,
  \textbf{cuando} actualiza campos válidos, \textbf{entonces} se
  conserva una nueva revisión y la anterior permanece consultable.
\item
  \textbf{AC-U2-05}: \textbf{Dado que} dos clientes parten de la misma
  revisión, \textbf{cuando} ambos intentan modificar, \textbf{entonces}
  una modificación basada en revisión obsoleta no sobrescribe la otra;
  se exige recargar.
\item
  \textbf{AC-U2-06}: \textbf{Dado que} una carrera tiene historial,
  \textbf{cuando} se desactiva con motivo, \textbf{entonces} conserva su
  identidad y referencias históricas; no se realiza eliminación física.
\item
  \textbf{AC-U2-07}: \textbf{Dado que} una identidad carece de capacidad
  de administración, \textbf{cuando} intenta crear/modificar/desactivar,
  \textbf{entonces} el catálogo no se modifica.
\item
  \textbf{AC-U2-08}: \textbf{Dado que} una identidad lectora solo tiene
  carrera A autorizada, \textbf{cuando} consulta listado, total o
  detalle de B, \textbf{entonces} no recibe información de B ni totales
  ajenos.
\end{enumerate}

\paragraph{User Story 3 - Preparar un proceso con modelo publicado
(Priority:
P2)}\label{user-story-3---preparar-un-proceso-con-modelo-publicado-priority-p2}

Como coordinador autorizado para gestionar procesos en una carrera,
quiero registrar carrera, período y versión publicada del modelo, para
preparar una autoevaluación identificable sin presumir que ya fue
autorizada institucionalmente.

\textbf{Why this priority}: utiliza el catálogo y modelo disponibles; es
posterior a las bases de acceso, pero puede validarse con fixtures sin
ejecutar altas nuevas ni activación.

\textbf{Independent Test}: utilizar una carrera activa y un modelo
publicado sintéticos preexistentes; crear el proceso y consultarlo como
el mismo actor, sin activar ni remitir evidencias.

\textbf{Acceptance Scenarios}:

\begin{enumerate}
\def\labelenumi{\arabic{enumi}.}
\tightlist
\item
  \textbf{AC-U3-01}: \textbf{Dado que} el actor está autorizado, la
  carrera/unidad están activas, el período pertenece a su vigencia y el
  modelo está publicado, \textbf{cuando} prepara el proceso,
  \textbf{entonces} se conserva con estado \path|PREPARATION|, carrera,
  período y versión exacta del modelo.
\item
  \textbf{AC-U3-02}: \textbf{Dado que} la carrera está desactivada o la
  unidad no está disponible, \textbf{cuando} intenta preparar un
  proceso, \textbf{entonces} se rechaza sin crear uno parcialmente
  válido.
\item
  \textbf{AC-U3-03}: \textbf{Dado que} el período termina antes de
  comenzar o excede la vigencia de la carrera, \textbf{cuando} intenta
  crear, \textbf{entonces} se rechaza sin guardar un proceso.
\item
  \textbf{AC-U3-04}: \textbf{Dado que} la versión del modelo no está
  publicada o no está disponible, \textbf{cuando} intenta crear,
  \textbf{entonces} se rechaza y no se sustituye automáticamente por
  otra versión.
\item
  \textbf{AC-U3-05}: \textbf{Dado que} existe un proceso en preparación
  de su ámbito, \textbf{cuando} lo consulta, \textbf{entonces} observa
  las referencias guardadas y ningún estado de activación o acreditación
  inferido.
\item
  \textbf{AC-U3-06}: \textbf{Dado que} no tiene capacidad/ámbito de
  gestión o consulta, \textbf{cuando} intenta crear o leer un proceso
  ajeno, \textbf{entonces} no produce cambios ni recibe información
  protegida.
\end{enumerate}

\paragraph{Edge Cases}\label{edge-cases}

\begin{itemize}
\tightlist
\item
  Diferencias de mayúsculas y espacios de un código no deben permitir
  duplicados conforme a la normalización definida.
\item
  \path|validUntil = null| representa vigencia sin término explícito; no
  equivale a cadena vacía ni a la cadena \path|"null"|.
\item
  Una carrera activa en estado puede no estar vigente en la fecha de
  consulta o requerir clasificación; no basta comprobar solo la palabra
  ACTIVE.
\item
  Crear proceso fija la referencia publicada; activar exige políticas
  adicionales excluidas aquí. No se mezcla publicación técnica con
  vigencia normativa institucional.
\item
  La revisión inicial de carrera es 1 y la del proceso actual es 0; no
  imponer una revisión inicial común a todos los recursos.
\item
  Paginación fuera de rango es inválida; un listado vacío autorizado no
  demuestra ausencia de registros en otros ámbitos.
\item
  Si falla una escritura crítica o su auditoría obligatoria, no debe
  presentarse como cambio completado sin evidencia consistente; el fallo
  requiere verificación transaccional específica.
\item
  La regla de cinco fallos ``consecutivos'' frente a acumulados en
  ventana de 15 min sigue pendiente; no se incluye una solución
  inventada.
\end{itemize}

\subsubsection{Requirements}\label{requirements}

\paragraph{Functional Requirements}\label{functional-requirements}

Los IDs \path|APE-FR-*| identifican la selección de esta práctica y no
reemplazan los FR canónicos. Se mantienen separados comportamiento
observable y detalle de implementación.

{\def\LTcaptype{none} % do not increment counter
\begin{longtable}[]{@{}
  >{\raggedright\arraybackslash}p{(\linewidth - 4\tabcolsep) * \real{0.3333}}
  >{\raggedright\arraybackslash}p{(\linewidth - 4\tabcolsep) * \real{0.3333}}
  >{\raggedright\arraybackslash}p{(\linewidth - 4\tabcolsep) * \real{0.3333}}@{}}
\toprule\noalign{}
\begin{minipage}[b]{\linewidth}\raggedright
ID local / origen
\end{minipage} & \begin{minipage}[b]{\linewidth}\raggedright
Requisito refinado
\end{minipage} & \begin{minipage}[b]{\linewidth}\raggedright
Verificación
\end{minipage} \\
\midrule\noalign{}
\endhead
\bottomrule\noalign{}
\endlastfoot
APE-FR-01 / FR-001,005 & El sistema deberá permitir iniciar/cerrar
sesión a identidades habilitadas y consultar sus funciones vigentes; al
cerrar o desactivar la identidad impedirá acceso posterior con esa
sesión. & AC-U1-01--04. \\
APE-FR-02 / FR-002,003 & Cada operación protegida deberá comprobar
capacidad y ámbito vigente; el acceso lector no concederá modificación
ni el privilegio técnico concederá competencia institucional. &
AC-U2-07/08, AC-U3-06. \\
APE-FR-03 / FR-043 & El administrador autorizado podrá registrar
unidades y carreras con código normalizado único y campos obligatorios;
una carrera requerirá unidad activa y vigencia válida. &
AC-U2-01--03. \\
APE-FR-04 / FR-043 & La consulta de catálogo deberá permitir identificar
carreras, unidades y vigencia y limitar listados, conteos y detalle al
ámbito efectivo del actor. & AC-U2-08; escenario lector y administrador
separados. \\
APE-FR-05 / FR-043 & Modificar o desactivar una carrera conservará
identidad, revisiones previas y relaciones históricas; desactivar
requerirá motivo y no eliminará físicamente la carrera. &
AC-U2-04/06. \\
APE-FR-06 / FR-010,012 & El actor autorizado podrá consultar e
identificar una versión publicada del modelo para utilizar su referencia
exacta al preparar el proceso. & AC-U3-01/04; fixture con versión
publicada. \\
APE-FR-07 / FR-011 & El actor con gestión vigente podrá crear un proceso
para carrera activa, período permitido y versión publicada; el resultado
quedará en preparación, sin activarse por el alta. & AC-U3-01--04. \\
APE-FR-08 / FR-002,011,012 & El actor autorizado podrá consultar el
proceso preparado y recuperar carrera, período, estado y versión del
modelo; una consulta ajena no revelará el recurso. & AC-U3-05/06. \\
APE-FR-09 / FR-006 & Los cambios críticos de catálogo y creación de
procesos deberán conservar actor, fecha, objeto, acción y resultado
mediante la auditoría definida para la operación. & Consulta de
historial y registros; fallo de auditoría por ensayar. \\
APE-FR-10 / FR-039 & Ante entrada inválida, falta de acceso o conflicto,
el sistema rechazará la operación con código y descripción estables y
correlación, sin exponer secretos ni datos ajenos. & Casos negativos de
US1--US3 y contratos CP-01--09. \\
\end{longtable}
}

\paragraph{Non-functional
Requirements}\label{non-functional-requirements}

{\def\LTcaptype{none} % do not increment counter
\begin{longtable}[]{@{}
  >{\raggedright\arraybackslash}p{(\linewidth - 4\tabcolsep) * \real{0.3333}}
  >{\raggedright\arraybackslash}p{(\linewidth - 4\tabcolsep) * \real{0.3333}}
  >{\raggedright\arraybackslash}p{(\linewidth - 4\tabcolsep) * \real{0.3333}}@{}}
\toprule\noalign{}
\begin{minipage}[b]{\linewidth}\raggedright
ID local / origen
\end{minipage} & \begin{minipage}[b]{\linewidth}\raggedright
Atributo y requisito verificable
\end{minipage} & \begin{minipage}[b]{\linewidth}\raggedright
Método de comprobación / límite
\end{minipage} \\
\midrule\noalign{}
\endhead
\bottomrule\noalign{}
\endlastfoot
APE-NFR-01 / RNF-001--002,005 & Credenciales con hash adaptativo y sal;
sesiones expiran tras 30 min de inactividad; respuestas y registros no
expondrán secretos. & Inspección de almacenamiento y pruebas con reloj;
no equivale a auditoría completa. \\
APE-NFR-02 / RNF-006 & p95 de resultado visible de las operaciones
interactivas seleccionadas ≤ 2 s bajo carga piloto, excluyendo archivos
e informes. & Fijar entorno y operaciones; la campaña histórica de
consultas no demuestra este subconjunto completo. \\
APE-NFR-03 / RNF-009 & Fallos de canales complementarios ajenos al
incremento no deberán impedir acceso, catálogo ni preparación. & Repetir
recorridos con canal indisponible; no se inventa porcentaje de
disponibilidad. \\
APE-NFR-04 / RNF-011 & Las modificaciones deberán detectar revisión
obsoleta y rechazarla sin pérdida silenciosa de cambios. & Dos clientes
parten de la misma revisión; uno debe recargar. \\
APE-NFR-05 / RNF-012--013 & Eventos de auditoría serán inmutables para
usuarios ordinarios y las solicitudes tendrán correlación para
investigar fallos sin datos sensibles. & Intento de modificación y
consulta autorizada; comprobar correlación. \\
APE-NFR-06 / RNF-014 & Los recorridos seleccionados deberán completarse
con teclado, foco visible y controles etiquetados, y verificarse contra
criterios aplicables WCAG 2.2 AA. & Revisión automática/manual; no se
declara conformidad completa por disponer de etiquetas. \\
APE-NFR-07 / RNF-017--018 & Las reglas de estados y autorización
permanecerán separadas de presentación/persistencia y tendrán pruebas
reproducibles con datos sintéticos. & Revisión de dependencias y
pruebas; reconocer límites de acoplamiento SQL. \\
APE-NFR-08 / contratos y FR-039 & La especificación de los intercambios
del incremento identificará recurso, operación, datos, resultados y
errores para que un cliente interprete su representación. & Comparar
CP-01--09 con controladores; refinamiento académico de
interoperabilidad, pendiente de ratificar como RNF independiente. \\
APE-NFR-09 / SC-007 & Las tareas de acceso, alta y preparación deberán
tener protocolo de usabilidad con participantes representativos y datos
sintéticos; el criterio global del proyecto es 19 de 20 participantes al
primer intento. & La aplicación de ese umbral al recorte U1 requiere
CL-APE-05; no se declara medición. \\
APE-NFR-10 / RNF-019 & Las demostraciones y pruebas del incremento
utilizarán identidades y datos sintéticos sin requerir información
institucional real. & Fixtures identificados y revisión de ejemplos; no
publicar tokens ni contraseñas reales. \\
\end{longtable}
}

\paragraph{Key Entities}\label{key-entities}

\begin{itemize}
\tightlist
\item
  \textbf{Identidad}: persona/cuenta, estado habilitado y referencia de
  autenticación; la desactivación conserva autoría.
\item
  \textbf{Asignación}: función/capacidad, ámbito y vigencia;
  preexistente en este recorte.
\item
  \textbf{Unidad académica}: identificador, código, nombre y estado.
\item
  \textbf{Carrera}: identificador, código, nombre, unidad, vigencia,
  estado y revisiones históricas.
\item
  \textbf{Versión de modelo}: identidad y número de una versión
  publicada; no se edita en este incremento.
\item
  \textbf{Proceso}: identidad, código, nombre, carrera, período, versión
  publicada y estado de preparación.
\item
  \textbf{Evento de auditoría}: actor, acción, objeto, fecha, resultado
  y correlación.
\end{itemize}

\subsubsection{Success Criteria}\label{success-criteria}

\paragraph{Measurable Outcomes}\label{measurable-outcomes}

\begin{itemize}
\tightlist
\item
  \textbf{APE-SC-01}: todos los escenarios funcionales AC-U1-01--04,
  AC-U2-01--08 y AC-U3-01--06 deberán tener resultado registrado antes
  de aceptar el incremento; esta entrega solo los especifica.
\item
  \textbf{APE-SC-02}: en los casos definidos de código normalizado
  duplicado se creará cero carreras adicionales.
\item
  \textbf{APE-SC-03}: las modificaciones exitosas conservarán una
  revisión anterior recuperable; la modificación obsoleta causará cero
  sobrescrituras silenciosas.
\item
  \textbf{APE-SC-04}: todos los procesos creados en el recorrido de alta
  quedarán en preparación y conservarán carrera, período y versión
  seleccionada; cero activaciones implícitas.
\item
  \textbf{APE-SC-05}: todos los requisitos locales tendrán origen
  canónico, historia/escenario y operación preliminar o método de
  verificación identificados.
\item
  \textbf{APE-SC-06}: cada decisión de clarificación resuelta tendrá
  persona, fecha, respuesta y respaldo; las no resueltas permanecerán
  pendientes, sin respuestas supuestas.
\item
  \textbf{APE-SC-07}: el objetivo p95 y los criterios de
  accesibilidad/usabilidad tendrán protocolo definido y resultado
  acotado; no se aceptará ``cumplido'' sin medición aplicable.
\end{itemize}

\subsubsection{Assumptions}\label{assumptions}

\begin{itemize}
\tightlist
\item
  Se utilizan usuarios y asignaciones sintéticos ya habilitados; no se
  exige crear una nueva autoridad.
\item
  Existe una versión publicada de demostración; no se presume vigencia
  normativa institucional.
\item
  El cliente dispone de conexión y entorno documentado; no se exige modo
  sin conexión.
\item
  El recorte U1 es propuesta de esta preparación, pendiente del docente;
  no es una decisión de una clase cuya evidencia no se aportó.
\item
  Los datos de ACD2 se usan como insumo técnico; las entregas
  ACD1/AA1/APE1 y su retroalimentación real deben incorporarse si el
  equipo las tiene.
\item
  Los contratos se recuperan del código operativo como punto de partida,
  pero no se convierten en OpenAPI definitivo ni certifican todos sus
  permisos.
\end{itemize}

\subsubsection{Clarifications}\label{clarifications}

No se ha efectuado una reunión con el docente en esta preparación. El
registro distingue decisión editorial adoptada de aprobación externa.

{\def\LTcaptype{none} % do not increment counter
\begin{longtable}[]{@{}
  >{\raggedright\arraybackslash}p{(\linewidth - 6\tabcolsep) * \real{0.2500}}
  >{\raggedright\arraybackslash}p{(\linewidth - 6\tabcolsep) * \real{0.2500}}
  >{\raggedright\arraybackslash}p{(\linewidth - 6\tabcolsep) * \real{0.2500}}
  >{\raggedright\arraybackslash}p{(\linewidth - 6\tabcolsep) * \real{0.2500}}@{}}
\toprule\noalign{}
\begin{minipage}[b]{\linewidth}\raggedright
ID
\end{minipage} & \begin{minipage}[b]{\linewidth}\raggedright
Pregunta / decisión
\end{minipage} & \begin{minipage}[b]{\linewidth}\raggedright
Estado y tratamiento
\end{minipage} & \begin{minipage}[b]{\linewidth}\raggedright
Validación requerida
\end{minipage} \\
\midrule\noalign{}
\endhead
\bottomrule\noalign{}
\endlastfoot
DEC-APE-01 & ¿Cómo se nombrará la entrega? & Adoptada documentalmente:
\path|specification.md| en esta carpeta académica; se conservan
\path|spec.md| canónicos. & Verificar que el docente acepta el enlace
como producto de práctica. \\
DEC-APE-02 & ¿Qué asignatura y datos se usan? & Se usan los de la guía
actual: Desarrollo Basado en Plataformas, 5 A, Unidad 1 y docente Edison
Leonardo Coronel Romero. & No trasladar la asignatura SDD de ACD2 como
nombre de esta guía. \\
CL-APE-01 & ¿Es suficiente acceso, catálogo y preparación sin activación
para Unidad 1? & Pendiente; se propone por viabilidad y dependencia de
autoridades no validadas. & Docente; respuesta, fecha y respaldo por
registrar. \\
CL-APE-02 & ¿Qué actores/asignaciones están institucionalmente
autorizados? & Pendiente CL-001/002/015; usar política sintética y no
activar institucionalmente. & Autoridad competente y docente según
alcance. \\
CL-APE-03 & ¿Se consideran cinco fallos consecutivos o acumulados en 15
min? & Pendiente: RNF-003 y SQL difieren; no cambiar política por
documentación. & Equipo/docente y responsable de seguridad. \\
CL-APE-04 & ¿Qué evidencia institucional sustenta el problema? &
Pendiente: no se aportaron entrevistas o actas. & Equipo; incorporar
fuente real sin inventar cifras. \\
CL-APE-05 & ¿Qué carga, entorno y tareas se medirán para RNF y
usabilidad? & Pendiente; conservar umbrales existentes y delimitar
protocolos. & Equipo/docente. \\
CL-APE-06 & ¿Qué correcciones indicó el docente al revisar esta
especificación? & Pendiente; no existe una observación de sesión
aportada. & Registrar texto, persona, fecha, IDs y cambio aceptado. \\
\end{longtable}
}

\subsubsection{Contratos preliminares}\label{contratos-preliminares}

Las fichas \href{CONTRATOS_PRELIMINARES.md}{CP-01--CP-09} integran este
documento por referencia. Identifican requisitos, recurso, operación,
entrada, resultado y errores. Su estado es preliminar académico,
contrastado con código; no se anuncia una API nueva ni se modifican
contratos productivos.

{\def\LTcaptype{none} % do not increment counter
\begin{longtable}[]{@{}
  >{\raggedright\arraybackslash}p{(\linewidth - 4\tabcolsep) * \real{0.3333}}
  >{\raggedright\arraybackslash}p{(\linewidth - 4\tabcolsep) * \real{0.3333}}
  >{\raggedright\arraybackslash}p{(\linewidth - 4\tabcolsep) * \real{0.3333}}@{}}
\toprule\noalign{}
\begin{minipage}[b]{\linewidth}\raggedright
Contrato
\end{minipage} & \begin{minipage}[b]{\linewidth}\raggedright
Operación / recurso
\end{minipage} & \begin{minipage}[b]{\linewidth}\raggedright
Requisitos
\end{minipage} \\
\midrule\noalign{}
\endhead
\bottomrule\noalign{}
\endlastfoot
CP-01 & Iniciar/cerrar sesión y consultar identidad. & APE-FR-01/02/10;
APE-NFR-01. \\
CP-02 & Consultar/registrar unidad académica. & APE-FR-03/04/09/10. \\
CP-03 & Registrar carrera. & APE-FR-03/09/10. \\
CP-04 & Consultar catálogo, detalle e historial. & APE-FR-04/05/10. \\
CP-05 & Modificar carrera con precondición. & APE-FR-05/09/10;
APE-NFR-04. \\
CP-06 & Desactivar carrera con motivo. & APE-FR-05/09/10. \\
CP-07 & Consultar modelo/versiones existentes. & APE-FR-06. \\
CP-08 & Preparar proceso. & APE-FR-07/09/10. \\
CP-09 & Consultar proceso preparado. & APE-FR-08/10. \\
\end{longtable}
}

\subsubsection{Trazabilidad inicial y verificación
posterior}\label{trazabilidad-inicial-y-verificaciuxf3n-posterior}

Las rutas son relativas a la raíz del repositorio; los nombres de
pruebas se citan como existentes, no como resultados nuevos.

{\def\LTcaptype{none} % do not increment counter
\begin{longtable}[]{@{}
  >{\raggedright\arraybackslash}p{(\linewidth - 6\tabcolsep) * \real{0.2500}}
  >{\raggedright\arraybackslash}p{(\linewidth - 6\tabcolsep) * \real{0.2500}}
  >{\raggedright\arraybackslash}p{(\linewidth - 6\tabcolsep) * \real{0.2500}}
  >{\raggedright\arraybackslash}p{(\linewidth - 6\tabcolsep) * \real{0.2500}}@{}}
\toprule\noalign{}
\begin{minipage}[b]{\linewidth}\raggedright
Necesidad
\end{minipage} & \begin{minipage}[b]{\linewidth}\raggedright
Requisitos / historia
\end{minipage} & \begin{minipage}[b]{\linewidth}\raggedright
Escenarios
\end{minipage} & \begin{minipage}[b]{\linewidth}\raggedright
Código y prueba candidata
\end{minipage} \\
\midrule\noalign{}
\endhead
\bottomrule\noalign{}
\endlastfoot
N-01 & APE-FR-01/02; US1 & AC-U1-01--04, AC-U2-07/08, AC-U3-06 &
\path|backend/src/main/java/ec/edu/unl/sec/identityaccess/application/SessionService.java|;
\path|backend/src/test/java/ec/edu/unl/sec/identityaccess/application/SessionServiceTest.java|;
\path|ScopeAuthorizationServiceTest.java| del mismo directorio. \\
N-02 & APE-FR-03--05; US2 & AC-U2-01--08 &
\path|backend/src/main/java/ec/edu/unl/sec/selfassessment/application/CareerAdministrationService.java|;
\path|backend/src/test/java/ec/edu/unl/sec/selfassessment/application/CareerAdministrationServiceTest.java|;
\path|tests/acceptance/career-catalog.spec.ts|. \\
N-03 & APE-FR-06--08; US3 & AC-U3-01--06 &
\path|backend/src/main/java/ec/edu/unl/sec/selfassessment/application/SelfAssessmentService.java|;
\path|backend/src/test/java/ec/edu/unl/sec/selfassessment/api/SelfAssessmentControllerTest.java|;
\path|tests/acceptance/complete-workflows.spec.ts|, subconjunto P05. \\
N-04 & APE-FR-09/10; APE-NFR-04/05 & AC-U2-04/05, negativos US1--US3 &
\path|backend/src/main/java/ec/edu/unl/sec/shared/api/EntityTagSupport.java|;
\path|backend/src/test/java/ec/edu/unl/sec/selfassessment/application/SelfAssessmentAuditServiceTest.java|. \\
\end{longtable}
}

La relación posterior será necesidad → requisito → historia → AC →
prueba identificada → ejecución con versión/entorno/resultado. Los RNF
requieren su verificación adicional; un endpoint presente o una prueba
con nombre parecido no acredita cumplimiento.

\subsubsection{Revisión de calidad y
cambios}\label{revisiuxf3n-de-calidad-y-cambios}

{\def\LTcaptype{none} % do not increment counter
\begin{longtable}[]{@{}
  >{\raggedright\arraybackslash}p{(\linewidth - 4\tabcolsep) * \real{0.3333}}
  >{\raggedright\arraybackslash}p{(\linewidth - 4\tabcolsep) * \real{0.3333}}
  >{\raggedright\arraybackslash}p{(\linewidth - 4\tabcolsep) * \real{0.3333}}@{}}
\toprule\noalign{}
\begin{minipage}[b]{\linewidth}\raggedright
Problema detectado
\end{minipage} & \begin{minipage}[b]{\linewidth}\raggedright
Refinamiento realizado
\end{minipage} & \begin{minipage}[b]{\linewidth}\raggedright
Motivo verificable
\end{minipage} \\
\midrule\noalign{}
\endhead
\bottomrule\noalign{}
\endlastfoot
``Gestionar carreras'' engloba alta, consulta y cambios. &
APE-FR-03/04/05 y AC-U2 separados. & Cada comportamiento tiene entrada,
resultado y rechazo identificables. \\
``Crear proceso'' puede interpretarse como activarlo. & APE-FR-07 exige
PREPARATION; activación fuera de recorte. & Cero activaciones implícitas
en APE-SC-04. \\
``Acceso seguro'' puede significar solo login. & Sesión, permiso, ámbito
y estados diferenciados. & Casos de sesión válida sin permiso y consulta
ajena. \\
``Rendimiento adecuado'' no define aprobación. & APE-NFR-02 mantiene p95
≤ 2 s y exige protocolo. & Medición reproducible; entorno aún
pendiente. \\
``Datos de demostración'' puede confundirse con decisiones oficiales. &
Supuestos y CL-APE-02/04 explícitos. & No se atribuye vigencia ni
competencia por fixture. \\
\end{longtable}
}

El documento es claro y pertinente al recorte propuesto; la
verificabilidad se expresa con AC e indicadores. La consistencia final
de alcance y políticas se mantiene pendiente de CL-APE. Las versiones
reales de trabajo y su diferencia se conservan en
\url{HISTORIAL_CAMBIOS.md}.

\subsection{Anexo B. Contratos
preliminares}\label{anexo-b.-contratos-preliminares}

\textbf{Estado:} fichas académicas preliminares, recuperadas del código
existente. No son un OpenAPI definitivo, una API nueva ni evidencia de
ejecución. Complementan \url{specification.md}.

\textbf{Base:} commit \path|faf868ea|; prefijo HTTP observado
\path|/api/v1|. Los códigos se obtuvieron de controladores, servicios y
manejadores de errores, no únicamente del OpenAPI, que presenta
divergencias documentadas en la auditoría.

\subsubsection{Convenciones y fuentes}\label{convenciones-y-fuentes}

\begin{itemize}
\tightlist
\item
  Solicitudes/respuestas de negocio en JSON; errores estructurados como
  Problem Details con \path|code|, \path|detail|, \path|correlationId| y
  campos HTTP relacionados.
\item
  Autenticación existente: cookie \path|SEC_SESSION|. Las mutaciones
  autenticadas requieren \path|X-CSRF-TOKEN|; no se supone JWT/Bearer.
\item
  Consultas de recurso revisable entregan ETag cuando el controlador lo
  implementa. Cambios existentes de carrera requieren \path|If-Match|
  fuerte, por ejemplo \path|"1"|.
\item
  Los nombres y DTO descritos reflejan el código actual. Los futuros
  cambios de esquema, versionado o política deben confirmarse; no se
  corrigen en esta práctica.
\item
  Los campos ID, autor, fecha y revisión los asigna el servidor; no se
  aceptan como decisiones de autoridad enviadas por el cliente.
\item
  Un ejemplo de resultado parcial en una ficha no reemplaza el DTO
  completo ni su esquema.
\end{itemize}

Rutas de fuente, relativas a la raíz:
{\def\LTcaptype{none} % do not increment counter
\begin{longtable}[]{@{}
  >{\raggedright\arraybackslash}p{(\linewidth - 2\tabcolsep) * \real{0.1500}}
  >{\raggedright\arraybackslash}p{(\linewidth - 2\tabcolsep) * \real{0.8500}}@{}}
\toprule\noalign{}
\begin{minipage}[b]{\linewidth}\raggedright
Referencia
\end{minipage} & \begin{minipage}[b]{\linewidth}\raggedright
Archivo
\end{minipage} \\
\midrule\noalign{}
\endhead
\bottomrule\noalign{}
\endlastfoot
API-S &
\path|backend/src/main/java/ec/edu/unl/sec/identityaccess/api/SessionController.java| \\
API-C &
\path|backend/src/main/java/ec/edu/unl/sec/selfassessment/api/CareerController.java| \\
API-M &
\path|backend/src/main/java/ec/edu/unl/sec/evaluationmodel/api/EvaluationModelController.java| \\
API-P &
\path|backend/src/main/java/ec/edu/unl/sec/selfassessment/api/SelfAssessmentController.java| \\
SVC-C &
\path|backend/src/main/java/ec/edu/unl/sec/selfassessment/application/CareerAdministrationService.java| \\
SVC-P &
\path|backend/src/main/java/ec/edu/unl/sec/selfassessment/application/SelfAssessmentService.java| \\
ERR-C &
\path|backend/src/main/java/ec/edu/unl/sec/selfassessment/api/CareerCatalogProblemHandler.java| \\
ERR-P &
\path|backend/src/main/java/ec/edu/unl/sec/selfassessment/api/ProcessExceptionHandler.java| \\
ERR-G &
\path|backend/src/main/java/ec/edu/unl/sec/shared/api/ApiProblemHandler.java| \\
REV &
\path|backend/src/main/java/ec/edu/unl/sec/shared/api/EntityTagSupport.java| \\
\end{longtable}
}

\subsubsection{CP-01 --- Sesión e
identidad}\label{cp-01-sesiuxf3n-e-identidad}

\begin{itemize}
\tightlist
\item
  Necesidad/requisitos: N-01; APE-FR-01/02/10, APE-NFR-01.
\item
  Operaciones: POST \path|/sessions|, GET \path|/me|, DELETE
  \path|/sessions|.
\item
  Entrada de login: \path|username| (3--100 caracteres), \path|password|
  (12--256), conforme al DTO. No publicar valores reales.
\item
  Resultado: POST devuelve 201, cookie y
  \path|{userId, expiresAt, csrfToken}|; GET 200 con
  identidad/funciones; DELETE 204 y cookie expirada, sin cuerpo.
\item
  Permisos: login público; consulta/cierre requieren sesión vigente.
  Logout autenticado exige CSRF.
\item
  Errores relevantes: 400 entrada inválida; 401 sesión/credenciales
  inválidas; 403 CSRF; 429 limitación de intentos. Los errores no
  confirman existencia de la cuenta.
\item
  Escenarios: AC-U1-01--04. Fuente: API-S y componentes de sesión.
\item
  Pendiente: resolver diferencia de consecutividad del umbral antes de
  declarar cumplimiento de RNF-003.
\end{itemize}

\subsubsection{CP-02 --- Unidad académica
disponible}\label{cp-02-unidad-acaduxe9mica-disponible}

\begin{itemize}
\tightlist
\item
  Necesidad/requisitos: N-02; APE-FR-03/04/09/10.
\item
  Operaciones: GET/POST \path|/academic-units|; GET
  \path|/academic-units/{id}|.
\item
  Entrada GET de colección: \path|page| (predeterminado 0), \path|size|
  (20) y \path|q| (texto). No se inventa filtro \path|status| para este
  listado.
\item
  Entrada POST: \path|{code, name}|. Requiere sesión, CSRF y capacidad
  de administración de catálogo.
\item
  Resultado: GET 200 con página o unidad; POST 201 con unidad, Location
  y ETag. Revisión inicial de unidad 1.
\item
  Errores: 400 datos/paginación inválidos; 401 sin sesión; 403 sin
  capacidad; 404 detalle no disponible/visible; 409 código duplicado.
\item
  Fuentes: API-C, SVC-C y ERR-C.
\item
  Nota: el alta de unidad puede evitarse en una prueba independiente
  mediante fixture activo. No se exige implementar una entidad nueva.
\end{itemize}

\subsubsection{CP-03 --- Registrar
carrera}\label{cp-03-registrar-carrera}

\begin{itemize}
\tightlist
\item
  Necesidad/requisitos: N-02; APE-FR-03/09/10.
\item
  Operación: POST \path|/careers|.
\item
  Entrada: \path|code| obligatorio hasta 100 caracteres; \path|name|
  obligatorio hasta 300; \path|academicUnitId| UUID obligatorio;
  \path|validFrom| fecha obligatoria; \path|validUntil| fecha o
  \path|null|.
\item
  Precondición: administrador autorizado, unidad activa, código
  normalizado único y vigencia coherente.
\item
  Resultado: 201 con Career; Location \path|/api/v1/careers/{id}|; ETag
  \path|"1"|; estado ACTIVE y revisión 1. Atributos de negocio se
  serializan dentro de \path|data|; el cuerpo de entrada es plano.
\item
  Errores: 400 campos/JSON inválidos; 401 sin sesión; 403
  capacidad/CSRF; 409 DUPLICATE\_CODE; 422 UNIT\_UNAVAILABLE cuando la
  unidad no está disponible.
\item
  Escenarios: AC-U2-01--03/07. Fuentes: API-C, SVC-C y
  \path|selfassessment/domain/Career.java|.
\end{itemize}

Ejemplo sintético de entrada, condicionado a que la unidad exista y esté
activa:

\begin{Verbatim}[breaklines=true,breakanywhere=true,fontsize=\footnotesize]
{
  "code": "U1-DEMO",
  "name": "Carrera sintética Unidad 1",
  "academicUnitId": "11111111-1111-4111-8111-111111111111",
  "validFrom": "2026-01-01",
  "validUntil": null
}
\end{Verbatim}

\subsubsection{CP-04 --- Consultar catálogo, detalle e
historial}\label{cp-04-consultar-catuxe1logo-detalle-e-historial}

\begin{itemize}
\tightlist
\item
  Necesidad/requisitos: N-02/N-04; APE-FR-04/05/10.
\item
  Operaciones: GET \path|/careers|; GET \path|/careers/{id}|; GET
  \path|/careers/{id}/history|.
\item
  Consulta de catálogo: \path|page|, \path|size|, \path|q|,
  \path|status|, \path|academicUnitId|. Estados del catálogo:
  ACTIVE/DISABLED.
\item
  Consulta de historial: \path|page|, \path|size|. El detalle usa ID de
  ruta, no cuerpo de solicitud.
\item
  Resultado: 200; listado/historial paginados según CatalogPage, detalle
  Career con ETag. Los totales deben corresponder al ámbito filtrado.
\item
  Permiso: capacidad de lectura vigente o administración según la
  política actual; no confundir un administrador global autorizado con
  un lector acotado a carrera.
\item
  Errores: 400 paginación/filtros inválidos; 401 sesión; 403 falta de
  capacidad de colección; 404 recurso ajeno/desconocido según servicio.
\item
  Escenarios: AC-U2-04/08. Fuentes: API-C y SVC-C.
\item
  Pendiente: revisar política completa de alcance administrativo
  institucional; el contrato preliminar no prueba todos sus controles.
\end{itemize}

\subsubsection{CP-05 --- Modificar carrera sin sobrescribir
cambios}\label{cp-05-modificar-carrera-sin-sobrescribir-cambios}

\begin{itemize}
\tightlist
\item
  Necesidad/requisitos: N-02/N-04; APE-FR-05/09/10 y APE-NFR-04.
\item
  Operación: PATCH \path|/careers/{id}|.
\item
  Entrada: ID de ruta; cabecera If-Match; mismo CareerInput completo de
  CP-03. Aunque sea PATCH, el controlador no acepta cualquier
  subconjunto de campos.
\item
  Precondición: administrador, sesión/CSRF y revisión vigente; unidad
  activa, código no duplicado y fechas válidas.
\item
  Resultado: 200 con Career, revisión incrementada y nuevo ETag;
  historial anterior conservado.
\item
  Errores: 400 datos/ETag malformados o débiles; 403 sin capacidad/CSRF;
  404 no encontrado; 409 código duplicado; 412 STALE\_REVISION; 422
  unidad no disponible; 428 falta If-Match.
\item
  Escenarios: AC-U2-04/05/07. Fuentes: API-C, SVC-C y REV.
\item
  Verificación: consultar revisión, modificar desde dos clientes y
  comprobar que la revisión obsoleta no sobrescribe el cambio exitoso.
  No se afirma ejecución.
\end{itemize}

\subsubsection{CP-06 --- Desactivar carrera con
motivo}\label{cp-06-desactivar-carrera-con-motivo}

\begin{itemize}
\tightlist
\item
  Necesidad/requisitos: N-02; APE-FR-05/09/10.
\item
  Operación: POST \path|/careers/{id}/state-decisions|.
\item
  Entrada: If-Match y
  \path|{status: "DISABLED", reason: "motivo de 10 a 4000 caracteres"}|.
\item
  Resultado: 201 con carrera, revisión incrementada y ETag; no se
  elimina su ID ni historial.
\item
  Errores: 400 validación de DTO; 403 sin capacidad/CSRF; 404 recurso
  ausente; 409 STATE\_UNCHANGED si no hay cambio; 412 revisión obsoleta;
  422 razón inválida cuando la detecta la regla; 428 sin precondición.
\item
  Escenario: AC-U2-06. Fuentes: API-C y SVC-C.
\item
  Límite: desactivar carrera no implica borrar ni cerrar automáticamente
  procesos anteriores. La reactivación existe en el servicio pero no
  amplía el recorte propuesto sin revisión del equipo.
\end{itemize}

\subsubsection{CP-07 --- Seleccionar una versión publicada
existente}\label{cp-07-seleccionar-una-versiuxf3n-publicada-existente}

\begin{itemize}
\tightlist
\item
  Necesidad/requisitos: N-03; APE-FR-06.
\item
  Operaciones: GET \path|/models|; GET
  \path|/models/{modelId}/versions|; GET \path|/model-versions/{id}|.
  Existen además variantes \path|/configured|, pero no se exige
  migrarlas en esta práctica.
\item
  Entrada: identidad autorizada, modelo o versión existente. Sin cuerpo
  en estas consultas.
\item
  Resultado: 200 con catálogo/lista/versión; el detalle entrega ETag.
  Identificar ID, número y estado de publicación.
\item
  Validación: el proceso solo puede prepararse con versión publicada
  disponible; no seleccionar automáticamente ``la última'' por nombre.
\item
  Fuentes: API-M, SVC-P y el puerto de catálogo publicado.
\item
  Límite: consultar o publicar técnicamente no demuestra aplicabilidad
  normativa para activación institucional. Publicación/activación quedan
  fuera del recorte U1.
\end{itemize}

\subsubsection{CP-08 --- Preparar proceso}\label{cp-08-preparar-proceso}

\begin{itemize}
\tightlist
\item
  Necesidad/requisitos: N-03; APE-FR-07/09/10.
\item
  Operación: POST \path|/processes|.
\item
  Entrada: \path|careerId|, \path|code| (hasta 100), \path|name| (hasta
  300), \path|periodStart|, \path|periodEnd|, \path|modelVersionId|,
  todos obligatorios. No se recibe un campo \path|state| para que el
  cliente fuerce activación.
\item
  Precondición: capacidad PROCESS\_MANAGE en carrera, identidad humana
  habilitada, carrera/unidad activas y período dentro de vigencia,
  versión publicada disponible.
\item
  Resultado: 201 con SelfAssessmentProcess; Location
  \path|/api/v1/processes/{id}|; estado PREPARATION; revisión inicial 0
  y ETag \path|"0"|. La referencia conserva \path|modelVersionId| y
  \path|modelVersionNumber|.
\item
  Errores relevantes: 400 DTO o MODEL\_VERSION\_NOT\_PUBLISHED; 401 sin
  sesión; 403 PROCESS\_SCOPE\_DENIED; 404 carrera inexistente/no visible
  según regla; 409 CAREER\_INACTIVE; 422 UNIT\_UNAVAILABLE o
  CAREER\_PERIOD\_MISMATCH. Diferenciar error de tipo/sintaxis de regla
  de período.
\item
  Escenarios: AC-U3-01--04/06. Fuentes: API-P, SVC-P, SVC-C y
  ERR-P/ERR-C.
\item
  Límite: aquí no se verifican todas las reglas de aplicabilidad
  normativa de activación; no declarar que el alta concede autorización
  institucional.
\end{itemize}

Ejemplo sintético de entrada, condicionado a fixtures existentes:

\begin{Verbatim}[breaklines=true,breakanywhere=true,fontsize=\footnotesize]
{
  "careerId": "22222222-2222-4222-8222-222222222222",
  "code": "PROCESO-U1-DEMO",
  "name": "Preparación sintética Unidad 1",
  "periodStart": "2026-01-01",
  "periodEnd": "2026-12-31",
  "modelVersionId": "33333333-3333-4333-8333-333333333333"
}
\end{Verbatim}

\subsubsection{CP-09 --- Consultar proceso
preparado}\label{cp-09-consultar-proceso-preparado}

\begin{itemize}
\tightlist
\item
  Necesidad/requisitos: N-03; APE-FR-08/10.
\item
  Operaciones: GET \path|/processes/{id}| y GET \path|/processes|.
\item
  Entrada: sesión y ID; el listado actual no declara paginación en su
  controlador. No inventar \path|page/size| como contrato implementado
  de procesos.
\item
  Resultado: 200 con proceso o lista filtrada; detalle con ETag.
  Permisos y proyección visible pueden diferir según
  funciones/indicadores. En preparación se consulta con un actor de
  gestión autorizado, no se presume acceso de un responsable sin
  indicadores asignados.
\item
  Errores: 401 sin sesión; 404 PROCESS\_NOT\_FOUND para recurso
  ausente/no visible según SVC-P.
\item
  Escenarios: AC-U3-05/06. Fuentes: API-P y SVC-P.
\end{itemize}

\subsubsection{Qué debe decidirse antes de convertir estas fichas en
contrato
definitivo}\label{quuxe9-debe-decidirse-antes-de-convertir-estas-fichas-en-contrato-definitivo}

{\def\LTcaptype{none} % do not increment counter
\begin{longtable}[]{@{}
  >{\raggedright\arraybackslash}p{(\linewidth - 4\tabcolsep) * \real{0.3333}}
  >{\raggedright\arraybackslash}p{(\linewidth - 4\tabcolsep) * \real{0.3333}}
  >{\raggedright\arraybackslash}p{(\linewidth - 4\tabcolsep) * \real{0.3333}}@{}}
\toprule\noalign{}
\begin{minipage}[b]{\linewidth}\raggedright
Punto
\end{minipage} & \begin{minipage}[b]{\linewidth}\raggedright
Evidencia disponible
\end{minipage} & \begin{minipage}[b]{\linewidth}\raggedright
Decisión pendiente
\end{minipage} \\
\midrule\noalign{}
\endhead
\bottomrule\noalign{}
\endlastfoot
Alcance de recursos de U1 & CP-01--09 sobre rutas existentes. &
Aceptación de CL-APE-01. \\
Nombres, esquemas y versión de API & Código v1 y OpenAPI común;
discrepancias conocidas. & Reconciliar sin romper clientes; no presentar
v2 como disponible. \\
Permisos administrativos & Política implementada y reglas
institucionales pendientes. & Confirmar competencia y alcance real. \\
Errores y paginación & 400/422 dependen de regla; catálogo paginado,
procesos no. & Documentar por operación; evitar una regla genérica
falsa. \\
Concurrencia & ETag/If-Match actuales. & Mantener revisiones de cada
agregado; no asumir todas empiezan en 1. \\
Validación de interfaces & Fuentes y pruebas existentes. & Ejecutar
posteriormente pruebas de esquema/HTTP en entorno sintético. \\
\end{longtable}
}

Los estándares HTTP y OpenAPI ayudan a describir operaciones y
representaciones, pero su existencia no obliga a que esta práctica
produzca ya un OpenAPI final. Referencias:
\href{https://www.rfc-editor.org/rfc/rfc9110.html}{RFC 9110} y
\href{https://spec.openapis.org/oas/latest.html}{OpenAPI}. No se cambia
la versión OpenAPI del repositorio por consultar el estándar actual.

\subsubsection{Desarrollo técnico disponible para revisión
personal}\label{desarrollo-tuxe9cnico-disponible-para-revisiuxf3n-personal}

En la preparación asistida se delimitó un incremento de acceso, catálogo
y preparación de proceso. Se separaron los requisitos de carrera en
alta, consulta y modificación histórica, y se hizo explícito que crear
proceso no significa activarlo. Estas tareas están desarrolladas en
\url{specification.md}, APE-FR-03--08 y escenarios AC-U2/AC-U3.

La diferencia verificable entre el borrador de esta preparación y la
especificación refinada está en
\url{evidencias/refinamiento-v01-v02.diff}. Es evidencia de los archivos
elaborados con apoyo de IA, no demuestra por sí sola la intervención
individual de un estudiante ni corresponde a cambios de una clase
anterior.

\section{Declaración de uso de Inteligencia
Artificial}\label{declaraciuxf3n-de-uso-de-inteligencia-artificial}

Se utilizó ChatGPT/Codex para inspeccionar documentación y código del
repositorio, seleccionar un incremento académico, redactar y refinar
requisitos, historias, criterios, contratos y respuestas de control, y
generar los archivos Markdown/LaTeX y su evidencia local. También se
consultaron fuentes técnicas públicas oficiales.

No se ejecutaron pruebas de negocio ni solicitudes HTTP, no se crearon
commits/PR ni se publicó la entrega. No se atribuye al docente una
revisión no realizada ni a estudiantes una contribución no identificada.
Los integrantes deben revisar el contenido, completar sus datos,
registrar su aporte y asumir responsabilidad por la versión presentada.
La revisión humana final continúa pendiente.

\end{document}
